\begin{lemma}
\label{lemma-derivative-zero-pth-power}
Let $k$ be a perfect field of characteristic $p > 0$.
Let $K/k$ be an extension.
Let $a \in K$. Then $\text{d}a = 0$ in $\Omega_{K/k}$
if and only if $a$ is a $p$th power.
\end{lemma}

\begin{proof}
By Lemma \ref{lemma-colimit-differentials} we see that there exists a subfield
$k \subset L \subset K$ such that $L/k$
is a finitely generated field extension and such that
$\text{d}a$ is zero in $\Omega_{L/k}$.
Hence we may assume that $K$ is a finitely generated field extension
of $k$.

\medskip\noindent
Choose a transcendence basis $x_1, \ldots, x_r \in K$
such that $K$ is finite separable over $k(x_1, \ldots, x_r)$.
This is possible by the definitions, see
Definitions \ref{definition-perfect} and
\ref{definition-separable-field-extension}.
We remark that the result holds for the purely transcendental
subfield $k(x_1, \ldots, x_r) \subset K$.
Namely,
$$
\Omega_{k(x_1, \ldots, x_r)/k} =
\bigoplus\nolimits_{i = 1}^r k(x_1, \ldots, x_r) \text{d}x_i
$$
and any rational function all of whose partial derivatives are zero
is a $p$th power. Moreover, we also have
$$
\Omega_{K/k} =
\bigoplus\nolimits_{i = 1}^r K\text{d}x_i
$$
since $k(x_1, \ldots, x_r) \subset K$ is finite separable
(computation omitted). Suppose $a \in K$ is an element such that
$\text{d}a = 0$ in the module of differentials. By our choice of $x_i$ we
see that the minimal polynomial $P(T) \in k(x_1, \ldots, x_r)[T]$
of $a$ is separable. Write
$$
P(T) = T^d + \sum\nolimits_{i = 1}^d a_i T^{d - i}
$$
and hence
$$
0 = \text{d}P(a) = \sum\nolimits_{i = 1}^d a^{d - i}\text{d}a_i
$$
in $\Omega_{K/k}$. By the description of
$\Omega_{K/k}$ above and the fact that $P$ was the minimal
polynomial of $a$, we see that this implies $\text{d}a_i = 0$.
Hence $a_i = b_i^p$ for each $i$. Therefore by
Fields, Lemma \ref{fields-lemma-pth-root}
we see that $a$ is a $p$th power.
\end{proof}
